% !TeX program = xelatex
\documentclass{article}
\usepackage{kotex}
\usepackage[a4paper, left=3cm, right=3cm, top=2.5cm, bottom=2.5cm]{geometry}
\usepackage{amsmath, amssymb}
\usepackage{tcolorbox}
\usepackage{array}
\usepackage{setspace}
\usepackage{xcolor}
\usepackage{titlesec}

\setstretch{1.3}
\titleformat{name=\section,numberless}
  {\normalfont\large\bfseries\color{blue}}{}
  {0pt}{}

\title{1.1 유클리드 공간}
\date{}
\author{}

\begin{document}
\maketitle

\begin{tcolorbox}[colback=blue!5, colframe=blue!60!black,
  title=제 1 장\quad 유클리드 기하학 --- 한눈에 보기,
  boxrule=0.5mm, arc=3mm]
\begin{tabular}{@{}ll@{}}
1.1  & 유클리드 공간 \\
1.2  & 등거리 변환과 합동 \\
1.3  & 평면에서의 반사 \\
1.4  & 공간에서의 반사 \\
1.5  & 평행 이동 \\
1.6  & 회전 \\
1.10 & 일반적인 등거리 변환 \\
\end{tabular}
\end{tcolorbox}

\section*{1.1 유클리드 공간}

\begin{tcolorbox}[
  colback=teal!4, colframe=teal!65!black,
  title=정의 \& 핵심 규칙,
  fonttitle=\bfseries, boxrule=0.5mm, arc=3mm]

\textbf{\textcolor{teal!65!black}{기호.}}\quad
$n$차원 유클리드 공간 $\mathbb{E}^n$은 실수의 $n$순서쌍 집합 $\mathbb{R}^n$과
본질적으로 동일하지만, 다음 차이가 있다.

\medskip
\begin{center}
\renewcommand{\arraystretch}{1.3}
\begin{tabular}{|c|c|}
\hline
$\mathbb{R}^n$ & $\mathbb{E}^n$ \\
\hline
좌표 구조와 원점이 고정됨 & 좌표·특수점 없는 순수 기하 공간 \\
수치 계산·좌표 사용 시 & 기하적 논의(좌표 없음) 시 \\
\hline
\end{tabular}
\end{center}

\noindent\textcolor{teal!50!black}{\rule{\linewidth}{0.4pt}}

\textbf{\textcolor{teal!65!black}{핵심 규칙 (동질성).}}\quad
유클리드 공간은 \textbf{동질적(homogeneous)}이다 --- 어느 장소도 다른 곳과 똑같아
보이므로, 기하적 대상을 \textbf{구부림·늘림·찢김 없이} 이동·비교할 수 있다.
이를 바탕으로 다음 측정 도구가 의미를 가진다.
\begin{itemize}\setlength\itemsep{0pt}
  \item \textbf{자} --- 두 점 사이의 거리 측정
  \item \textbf{각도기} --- 각의 크기 측정
  \item \textbf{회전 방향} --- 시계 방향/반시계 방향 구별
\end{itemize}

\noindent\textcolor{teal!50!black}{\rule{\linewidth}{0.4pt}}

\textbf{\textcolor{teal!65!black}{함수 표기.}}\quad
$f : \mathbb{E}^n \to \mathbb{E}^n$은 각 점 $P \in \mathbb{E}^n$을 새 위치
$f(P) \in \mathbb{E}^n$으로 옮긴다. 집합 $S$의 상(image)은
\[
  f(S) \;=\; \{\, f(P) \mid P \in S \,\}
\]

\end{tcolorbox}

\end{document}
